\documentclass[10pt,a4paper]{article}
\usepackage[utf8]{inputenc}
\usepackage[T1]{fontenc}
\usepackage[brazil]{babel}
\usepackage{mathptmx}
\usepackage[a4paper,margin=1.6cm]{geometry}
\usepackage{enumitem}
\setlength{\parindent}{0pt}
\setlength{\parskip}{0.35em}
\pagestyle{empty}

\begin{document}

{\large\bfseries Referências citadas e o que cada uma conversa com o trabalho}\\[0.2em]
{\small Classificação bayesiana dinâmica de proficiência em Biologia -- EST613}

\vspace{0.6em}

\textbf{Van Moere \& Hanlon (2020)} -- \textit{Language Testing}, 37(4), 482--502.\\
\textbf{O que fizeram:} combinaram, pelo teorema de Bayes, as notas anteriores (como priori) com a nota mais recente, considerando o erro de medida, para atualizar a habilidade de mais de 21 mil alunos que fizeram um teste de leitura de 3 a 5 vezes. As estimativas ficaram mais precisas que usar só a última nota.\\
\textbf{Com o trabalho:} sustenta o uso do \emph{histórico como priori}, atualizado a cada avaliação.

\textbf{Zitzmann et al. (2025)} -- \textit{Educational and Psychological Measurement}, 85(3), 483--506.\\
\textbf{O que fizeram:} propuseram avaliar o aluno combinando as respostas na prova com a crença prévia do avaliador e compararam, por simulação, com a teoria clássica dos testes. A abordagem bayesiana é melhor \emph{apenas quando o avaliador é pouco viesado} e a confiança na priori é moderada.\\
\textbf{Com o trabalho:} explica por que o M1 (histórico sem ajuste, viesado) fica perto da média simples e o M2 (que desconta escola e professor) dá o salto.

\textbf{Kaplan et al. (2023)} -- \textit{Psychometrika}, 88(1), 1--30.\\
\textbf{O que fizeram:} compararam formas de incorporar dados históricos em avaliações de larga escala (PISA). O \emph{empréstimo dinâmico}, em que o quanto se aproveita do passado é estimado, foi tão bom ou melhor que agregar tudo ou usar \emph{power priors}, sobretudo quando os dados históricos são heterogêneos.\\
\textbf{Com o trabalho:} o M2 estima o quanto confiar no histórico ($\sigma_h$, $\tau$), enquanto o M1 fixa esse valor.

\textbf{Kalman (1960)} -- \textit{Journal of Basic Engineering}, 82(1), 35--45.\\
\textbf{O que fez:} propôs o filtro que prevê o estado de um sistema, observa uma medida e corrige a previsão, ponderando pela incerteza de cada parte.\\
\textbf{Com o trabalho:} é a mecânica da atualização prova a prova (priori $\to$ posteriori $\to$ priori da próxima), com o ruído $\omega$ permitindo que o aluno mude.

\textbf{Dias et al. (2022)} -- \textit{Theoretical and Applied Genetics}, 135, 1385--1399.\\
\textbf{O que fizeram:} recomendaram cultivares a partir da probabilidade a posteriori de desempenho superior, calculada com amostras do MCMC.\\
\textbf{Com o trabalho:} mesma lógica da probabilidade de cada nível (proporção das amostras em cada faixa), usada para decidir a classificação.

\textbf{Gelman (2006)} -- \textit{Bayesian Analysis}, 1(3), 515--534.\\
\textbf{O que fez:} mostrou que, com poucos grupos, a priori do desvio-padrão entre grupos influencia o resultado e recomendou half-Cauchy/half-Normal no lugar da Inv-Gama vaga.\\
\textbf{Com o trabalho:} justifica a análise de sensibilidade de $\tau$ (half-Normal $\times$ half-Cauchy), já que há só 3 professores atuais e 5 escolas.

\textbf{Spiegelhalter et al. (2002)} -- \textit{JRSS B}, 64(4), 583--639.\\
\textbf{O que fez:} propôs o DIC, critério que equilibra ajuste e complexidade (número efetivo de parâmetros).\\
\textbf{Com o trabalho:} primeiro critério de comparação dos modelos, sem usar a verdade da simulação.

\textbf{Vehtari, Gelman \& Gabry (2017)} -- \textit{Statistics and Computing}, 27, 1413--1432.\\
\textbf{O que fizeram:} propuseram calcular a validação cruzada deixando um de fora (LOO) de forma aproximada (PSIS), com o diagnóstico Pareto-$k$ para indicar quando a aproximação não é confiável.\\
\textbf{Com o trabalho:} LOO por aluno (``o modelo prevê bem um aluno novo?''); o Pareto-$k$ justificou descartar o LOO por nota.

\textbf{Carpenter et al. (2017)} -- \textit{Journal of Statistical Software}, 76(1), 1--32.\\
\textbf{O que fizeram:} apresentaram o Stan, linguagem para modelos bayesianos com o amostrador NUTS (Hamiltonian Monte Carlo).\\
\textbf{Com o trabalho:} software usado para ajustar o M2.

\vfill
{\small\textbf{Em uma frase:} a ideia de histórico como priori vem de Van Moere, a mecânica de atualização vem de Kalman, e o porquê de corrigir escola e professor vem de Zitzmann.}

\end{document}
